\input{preamble}

\usepackage{komorebi-monograph}
\renewcommand{\KomorebiKind}{Lean and evidence guide}
\renewcommand{\KomorebiSubtitle}{Exact coordinates, singular packets,\\and stable reconstruction}
\renewcommand{\KomorebiShortTitle}{Paper V / Lean guide}
\renewcommand{\KomorebiCoverBlurb}{Read the exact declarations, follow the signal models they realize, and separate preserved kernel evidence from current source and computation checks.}
\KomorebiSetup
\setlength{\cftsubsecnumwidth}{4.4em}
\definecolor{teal}{HTML}{245D60}
\setmonofont{DejaVuSansMono}[Path=fonts/,UprightFont=*.ttf,BoldFont=*-Bold.ttf,ItalicFont=*-Oblique.ttf,BoldItalicFont=*-BoldOblique.ttf,Scale=.80]
% Collected core edition: typography only. Original scientific body is unchanged.
\input{publisher-layout.tex}
\CoreLayout
\colorlet{accent}{pubTeal}\colorlet{teal}{pubTeal}\colorlet{ink}{pubInk}\colorlet{muted}{pubMuted}\colorlet{light}{pubWash}\colorlet{pale}{pubWash}
\renewcommand{\KomorebiCover}{%
 \clearpage\hypersetup{pageanchor=false}\thispagestyle{empty}\begingroup\setlength{\parindent}{0pt}
 {\sffamily\fontsize{8.5}{11}\selectfont\color{pubTeal}\MakeUppercase{\KomorebiSeries}\par}
 \vspace{35pt}
 {\fontsize{28}{34}\selectfont\color{pubNavy}\raggedright\KomorebiCoverTitle\par}
 \vspace{16pt}
 {\sffamily\fontsize{11}{15.5}\selectfont\color{pubMuted}\raggedright\KomorebiSubtitle\par}
 \vspace{20pt}\noindent\textcolor{pubCopper}{\rule{42pt}{.8pt}}\par
 \vspace{20pt}
 {\fontsize{11.4}{16.2}\selectfont\raggedright\KomorebiCoverBlurb\par}
 \vspace{24pt}{\small\sffamily\color{pubMuted}\textbf{Edition locators.} The contents and live internal links use this typeset component. Fixed page locators quoted from earlier reader editions retain their original pagination. Traditional-paper locators refer to the supplied frozen reader: \par\smallskip\nolinkurl{frozen_readers/Paper_V_Traditional_Final.pdf}\par}
 \vfill
 {\fontsize{14}{18}\selectfont Komorebi P.\par}
 \vspace{7pt}{\sffamily\fontsize{9}{12}\selectfont\color{pubMuted}\KomorebiKind\quad\textbullet\quad\KomorebiEdition\par}
 \endgroup\clearpage\pagenumbering{roman}\hypersetup{pageanchor=true}\color{pubInk}\normalfont
}

\begin{document}
\KomorebiCover
\KomorebiOpeningNote{
\textbf{A guide to the actual formal development.}
20 scoped manuscript results. 225 theorem targets. 23 proof modules.
The eight tutorial chapters explain the route from literal samples to Hilbert operators, finite filters, packet synthesis, adaptive trees, and projective actions. The complete ledgers then give exact declarations and manuscript locators.

The full traditional edition also contains 63 broader prose proof records. Their separate inventory is a reading aid and adds no formal coverage. The two edge recovery conjectures remain open.

This edition preserves 52 exact Lean excerpts and reconciles source bytes, declaration locations, dependency pins, and the saved axiom audit. Kernel execution evidence comes from the recorded formal-core build. Its toolchain, source identities and replay instructions appear in the reproduction section. This reading guide does not claim a new kernel execution.}

\setcounter{tocdepth}{1}
\begingroup\small
\tableofcontents
\endgroup
\clearpage
\pagenumbering{arabic}
\section{How to read the formal companion}
\label{sec:reading}
Paper V places several kinds of mathematics in one signal processing language: coordinate partitions, bounded Hilbert space operators, finite filter banks, finite tree optimization, formal series, and projective actions. The formal development keeps their types and hypotheses distinct. Following those distinctions is the most useful way to read its proofs.

A pointwise partition identity needs little structure. An energy theorem needs a normed space with the intended norm. A transform theorem needs convergence. An optimal finite pruning needs a finite family and an objective. A fractional linear expression needs a chart, while its projective action survives at the chart's poles. Each transition has a named construction or theorem in the source.

\begin{keyidea}[title={Follow the model before the theorem name}]
Start with the ambient type, then locate the bridge to the intended signal model. A coordinate identity, a bounded operator theorem, and an analytic transform law have different inputs. The exact excerpt shows where those inputs enter the declaration.
\end{keyidea}

\subsection{What the completed contract contains}
The selected release, \code{1.1.0-expanded-lean-checked}, records twenty scoped theorem-like manuscript results with explicit proofs and 225 supporting theorem targets. The inherited namespace audit reports 827 declarations. These are different denominators: the 827 includes definitions and compiler-generated declarations, and the 225 includes supporting results rather than 225 independent new mathematical claims.

The formal core is the expanded manuscript preserved as Part I of the final traditional edition. Part II preserves the broader Annals development. Its additional proof records are prose mathematics. They do not enter the twenty-result formalization contract. Section~\ref{sec:extension} gives a separate crosswalk.

The two edge recovery conjectures remain open. Known-edge cancellation is established, measurable-mask perturbation is established, and finite adaptive selection is established within its hypotheses. Those results do not provide an estimator with the full statistical and approximation guarantees requested by the conjectures.

\subsection{Evidence available in this edition}
\begin{longtable}{>{\raggedright\arraybackslash}p{.24\textwidth}>{\raggedright\arraybackslash}p{.66\textwidth}}
\toprule Evidence & Meaning \\ \midrule\endhead
Source reconciliation & Current hashes, 225 declaration locations, 225 audit directives, and 225 inherited axiom-output records agree. \\ \addlinespace
Historical Lean execution & The September 6 UTC release records the pinned build and namespace audit passing. The raw logs are preserved. \\ \addlinespace
Manual correspondence & Manuscript statements have a reviewed, explicit relationship to formal declarations. Lean does not parse or certify the English PDF automatically. \\ \addlinespace
Exact model tests & Finite independent models exercise indexing, masks, trees, lifting, filters, and boundary cases. They are computational evidence on their tested ranges. \\ \addlinespace
Companion numerics & Signal experiments and figure data are empirical demonstrations. Their floating-point implementation is outside the Lean theorem claim. \\ \bottomrule
\end{longtable}
The guide's current build performs no Lean compilation. Reproduction commands appear in Section~\ref{sec:reproduce}, with the missing-tool case treated as a failure to establish a fresh kernel result. The final package's verification records distinguish current reruns from inherited evidence.

\subsection{Read a theorem as a typed contract}
The name after \code{theorem} is a reusable declaration. Parenthesized arguments are explicit inputs. Braced arguments are generally inferred. A hypothesis such as \code{(hN : 0 < N)} is an actual input to the theorem. A typeclass argument such as \code{[AddCommMonoid E]} specifies the algebraic laws available to its proof.

The expression after the colon is the proposition. The proof after \code{:= by} constructs evidence of that proposition. Common commands here have small, specific jobs: \code{ext} reduces equality to coordinates, \code{rw} rewrites using an equality, \code{simp} uses declared simplification rules, \code{ring} normalizes polynomial identities, \code{omega} handles integer arithmetic, and \code{linarith} combines linear inequalities. Their successful use belongs to the recorded compilation, not to a textual count of tactic names.

A definition marked \code{noncomputable} may use classical choice or real-number constructions that do not give executable numerical code. Its theorems can still be checked by the kernel. In particular, the pruning evaluator over exact real costs specifies a mathematical algorithm and a count of unit operations. Lean evaluation speed is outside that count.

\subsection{Four kinds of signal space}
\begin{longtable}{>{\raggedright\arraybackslash}p{.25\textwidth}>{\raggedright\arraybackslash}p{.29\textwidth}>{\raggedright\arraybackslash}p{.35\textwidth}}
\toprule Object & Lean representation & What it provides \\ \midrule\endhead
Arbitrary sequence & \code{ℤ → K} & Every integer sample exists. No norm or convergence is assumed. \\ \addlinespace
Finite filter & \code{ℤ →₀ K} & A finitely supported coefficient function. Convolution sums are finite. \\ \addlinespace
Square-summable signal & \code{lp (fun _ : I => ℂ) 2} & The joint $\ell^2$ norm, infinite norm-convergent expansions, bounded maps and adjoints. \\ \addlinespace
Formal series & \code{R⟦X⟧} & Coefficients with formal addition and Cauchy multiplication. Analytic convergence is not part of the type. \\ \bottomrule
\end{longtable}
The same informal word ``packet'' appears in all four settings. The guide always identifies the ambient type before using an identity across them.

\subsection{Recommended reading routes}
For the signal processing route, read integer addresses, the Hilbert bridge, FIR transfer, faithful filter banks, and Haar packets. For the adaptive route, read known edges, mask perturbations, complete prunings, the dynamic program, and its stability envelope. For the algebraic route, read formal cyclic products, homogeneous Smith actions, and CRT masks.

Every displayed code excerpt is copied from an immutable source declaration. The raw excerpt, its source line range, and its SHA256 hash accompany the PDF. Excerpts rely on the imports, namespaces, variables, and definitions in their source files. They are not advertised as independently compilable snippets. The complete project remains the authoritative executable object.

\Needspace{10\baselineskip}\section{Partitions, signed addresses, and the Hilbert bridge}
\subsection{A packet first means a coordinate mask}
Take an index set $I$, a label set $R$, and a phase map $\phi:I\to R$. The packet with label $r$ keeps coordinate $i$ exactly when $\phi(i)=r$:
\[
(P_r x)(i)=\begin{cases}x(i),&\phi(i)=r,\\0,&\phi(i)\ne r.\end{cases}
\]
The definition requires a zero in the value type and a decidable label equality. It does not require a Fourier transform, a measure, or an inner product.
\excerpt{Packets}{packet}
Applying the same mask twice changes nothing. Distinct labels cannot both be retained at the same coordinate. If the label type is finite and values can be added commutatively, the finite sum of packets returns the original value. At any one coordinate exactly one summand survives.
\excerpt{Packets}{packet_complete}
The proof uses \code{Finset.sum_eq_single} to isolate the surviving label. This simple mechanism is reused after stronger ambient structures are introduced.

\subsection{Signed polyphase coordinates are an equivalence}
For $N>0$, every integer has the unique decomposition
\[
n=Nq+r,\qquad 0\le r<N,\quad q\in\mathbb Z.
\]
The integer quotient is Euclidean: for $N=3$ and $n=-1$, the address is $(r,q)=(2,-1)$, not a negative residue. Lean represents the bounded residue by \code{Fin N}. Its stored value has a proof that it is smaller than $N$.
\excerpt{IntegerPolyphase}{address}
The two inverse laws are mathematical content. They ensure both that analyzing then synthesizing recovers a signal and that synthesizing then analyzing recovers an arbitrary channel family. Merely checking the first direction would not establish that all channel coordinates are reachable.

Writing $x_r[q]=x[Nq+r]$, the pointwise maps are
\[
Ax(r,q)=x[Nq+r],\qquad Sc(n)=c(n\bmod N,\lfloor n/N\rfloor).
\]
These are meaningful for any value type. Energy is a later statement.

\subsection{The two units of delay}
A delay by $d$ original samples is $(D_dx)[n]=x[n-d]$. A delay by $d$ inside each channel is an original delay by $Nd$. For a general original delay, the residue changes and a signed carry enters the channel coordinate:
\[
(A D_d x)_r[q]
=x_{(r-d)\bmod N}\!\left[q+\left\lfloor\frac{r-d}{N}\right\rfloor\right].
\]
\excerpt{IntegerPolyphase}{sample_delay}
For $N=2$ and a one-sample delay, the new even channel is the old odd channel delayed once, while the new odd channel is the old even channel at the same channel index:
\[
(A D_1x)_0[q]=x_1[q-1],\qquad (A D_1x)_1[q]=x_0[q].
\]
For a two-sample delay both channels are delayed once. This is why a filter-bank identity with channel delay $d$ reconstructs $x[n-Nd]$.

\subsection{Square summability is built into the type}
\excerpt{HilbertPackets}{signal}
The bounded projector is obtained by first proving that masking cannot increase the norm. The proof works coordinate by coordinate and invokes monotonicity of the $\ell^2$ norm.
\excerpt{HilbertPackets}{projection_norm}
The same mask is made into a linear map and then promoted to a continuous linear map with bound one.
\excerpt{HilbertPackets}{projector}
The notation \code{→L[ℂ]} records continuous complex linearity. Idempotence alone would not establish boundedness, symmetry, or orthogonality. The file proves those separately. With finitely many labels, the sum is the identity as a continuous linear operator, not just a formula on finite test vectors.

\subsection{The channel norm is the joint norm}
The Hilbert space analysis target is $\ell^2(\operatorname{Fin}N\times\mathbb Z)$:
\[
\|Ax\|^2=\sum_{r=0}^{N-1}\sum_{q\in\mathbb Z}|x[Nq+r]|^2=\|x\|^2.
\]
It is not an unqualified function space of channels carrying a supremum norm. The source packages the address equivalence into a complex linear isometric equivalence.
\excerpt{HilbertPackets}{polyphase}
Its inverse laws are inherited from that equivalence, and its energy law follows from norm preservation. The proof of reindexing uses summability over the bijective address map. This is the actual bridge from arbitrary pointwise coordinates to the Hilbert statement.

\subsection{Impulse expansion means norm convergence}
The impulse at $i$ with amplitude $x_i$ is represented by \code{lp.single 2 i (x i)}. The full theorem is a \code{HasSum} statement in the $\ell^2$ space.
\excerpt{HilbertPackets}{impulse}
Thus the series of impulses converges to $x$ in the ambient norm. The finite support version is an exact finite sum. The shift is an isometric equivalence with formula $S_dx[n]=x[n-d]$, and shifting the impulse at zero gives the impulse at $d$. Density of the impulse basis yields the dense shift-generated span.

\subsection{Refinement uses divisibility}
If $N\mid M$, reducing modulo $M$ and then modulo $N$ agrees with reducing directly modulo $N$. Consequently one coarse residue projector is the sum of the fine residue projectors lying above it.
\excerpt{HilbertPackets}{refinement}
For example, the even packet modulo $2$ is the sum of packets $0$, $2$, and $4$ modulo $6$. The theorem is an equality of bounded operators on the full complex $\ell^2$ space. The hypotheses $N>0$, $M>0$, and $N\mid M$ remain explicit. A relation between arbitrary unrelated moduli would require a different construction.

\Needspace{10\baselineskip}\section{FIR transfer and faithful perfect reconstruction}
\subsection{Finite support makes time domain transfer unconditional}
An FIR filter is a finitely supported function $h:\mathbb Z\to K$. Negative tap addresses are allowed, so the formal theorem includes advances as well as delays. For an arbitrary bilateral sequence,
\[
(h*x)[n]=\sum_{t\in\operatorname{supp}h}h[t]x[n-t].
\]
No convergence hypothesis is needed because this sum is finite.
\excerpt{FIRTransfer}{fir}
Fix output phase $k$ and input phase $s$. The tap carrying an input channel delay $m$ is
\[
t=Nm+k-s,\qquad h_{k,s}[m]=h[Nm+k-s].
\]
The address identity is
\[
(Nn+k)-(Nm+k-s)=N(n-m)+s.
\]
It says exactly which original tap contributes to which channel and delay. The source's \code{tapEquiv} proves the relevant reindexing is bijective, then partitions the finite support by input phase. It does not enlarge the finite sum to an arbitrary infinite series.
\excerpt{FIRTransfer}{fir_law}

\subsection{A two channel calculation}
\begin{microexample}[title={Three taps and their channel delays}]
Take $h[0]=a$, $h[1]=b$, $h[2]=c$, with all other taps zero. For $N=2$ the four channel kernels are
\[
\begin{aligned}
y_0[n]&=a x_0[n]+c x_0[n-1]+b x_1[n-1],\\
y_1[n]&=b x_0[n]+a x_1[n]+c x_1[n-1].
\end{aligned}
\]
The tap at one has different channel delays in the two off diagonal positions. In the transform convention used by the paper, the transfer matrix is
\[
H(z)=\begin{pmatrix}a+c z^{-1}&b z^{-1}\\b&a+c z^{-1}\end{pmatrix}.
\]
The time domain equations hold for every sequence. The transform equation $Y(z)=H(z)X(z)$ has the additional analytic hypotheses described next.

\end{microexample}

\subsection{A transform is a convergent weighted sum}
The source defines a totalized \code{tsum} expression, but the meaningful convergence predicate is \code{HasTransform}.
\excerpt{ZTransform}{has_transform}
It states that the bilateral weighted series $\sum_n x[n]z^{-n}$ has the specified sum. The delay theorem requires both a nonzero transform parameter and this convergence input.
\excerpt{ZTransform}{delay_transform}
Finitely many delayed inputs can then be combined using finite sums of \code{HasSum} proofs. The matrix version assumes weighted summability separately for every input channel.
\excerpt{ZTransform}{matrix_transform}
An arbitrary $\ell^2$ signal need not have an absolutely convergent bilateral transform at a chosen complex $z$. An $L^2$ Fourier boundary-value interpretation on the unit circle is a separate analytic interpretation. Off the unit circle, the displayed weighted summability conditions must be supplied. Neither the definition of \code{tsum} nor an FIR input filter removes them.

\subsection{Laurent polynomials act as finite filters}
The type \code{AddMonoidAlgebra ℂ ℤ} represents finite Laurent coefficients indexed by integer delays. Its multiplication is finite convolution. A monomial at index $d$ acts by $x[n]\mapsto x[n-d]$, with transform symbol $z^{-d}$. The formal action is an algebra homomorphism into complex linear endomorphisms of all bilateral sequences.
\excerpt{FilterBank}{act}
The source proves \code{act_is_fir}, making the abstract algebra homomorphism coincide with the literal finite filter sum. It then defines a matrix bank by summing the action of each matrix entry over the input channels. The theorem \code{bank_mul} identifies matrix multiplication with composition of banks in the correct order.

\subsection{Why the matrix criterion is an equivalence}
The reverse direction of a perfect-reconstruction criterion needs faithfulness. If two matrix banks have the same action on every sequence, why must their Laurent coefficients agree? The source feeds a unit impulse into one channel at a time. Reading output channel $k$ at sample $n$ recovers exactly the $n$th Laurent coefficient of the $(k,s)$ entry.
\excerpt{FilterBank}{faithful}
This eliminates a possible gap between an operator statement and a formal matrix identity. It relies on equality on every input, rather than agreement on a bounded collection of numerical examples.

\subsection{Gain and delay}
The target matrix is diagonal, with each diagonal entry the monomial of amplitude $c$ at delay $d$. The theorem gives the equivalence between matrix equality and reconstruction of every channel sequence.
\excerpt{FilterBank}{perfect}
The formula allows any complex $c$. To interpret the result as invertible reconstruction up to gain, one additionally requires $c\ne0$. At $c=0$, the theorem remains a valid zero-output identity, but it is not an invertibility statement.

Returning from channel coordinates to original coordinates yields the literal sample formula.
\excerpt{FilterBank}{sample_reconstruction}
This result is independent of numerical transform convergence. The transform representation of the bank is a further theorem with \code{HasTransform} hypotheses. The distinction is practical: an exact finite reconstruction identity can be checked algebraically even when a particular signal's off-circle transform does not exist.

\Needspace{10\baselineskip}\section{Haar, lifting, and finite depth packet synthesis}
\subsection{Interleaving is part of the theorem}
For a bilateral signal, write $e[n]=x[2n]$ and $o[n]=x[2n+1]$. The source first proves that the squared norms of these two channels sum to the original squared norm. It then forms the normalized Haar channels
\[
a[n]=\frac{e[n]+o[n]}{\sqrt2},\qquad d[n]=\frac{e[n]-o[n]}{\sqrt2}.
\]
The sample formula establishes that the abstract construction really acts on these literal interleaved samples.
\excerpt{HaarChannels}{haar_samples}
The energy identity combines the Hilbert space parallelogram law with the joint channel norm.
\excerpt{HaarChannels}{haar_energy}
This theorem holds for complex $\ell^2$ signals, not only finite real arrays. A finite numerical check can expose normalization mistakes but cannot replace the infinite-space proof.

\subsection{Lifting is a different normalization}
The lifting analysis sends $(e,o)$ to $(s,d)$ with
\[
d=o-e,\qquad s=e+\frac d2=\frac{e+o}{2}.
\]
The inverse is $e=s-d/2$, $o=d+e$. Its determinant is one. The theorem is stated over a characteristic zero field, which makes division by two meaningful in the required way.
\excerpt{Lifting}{lifting}
The lifting pair is invertible but is not the normalized Haar pair. For example, $(e,o)=(2,4)$ gives lifting coordinates $(s,d)=(3,2)$, whose squared Euclidean norm is $13$, whereas the input norm squared is $20$. Normalized Haar gives $(3\sqrt2,-\sqrt2)$ with squared norm $20$ under the source's high-channel sign convention. Invertibility and isometry are separate claims.

\subsection{Flat finite signals have both inverse laws}
The finite transform file builds a pairing equivalence between a flat vector of length $2N$ and $N$ pairs. It transports lifting through that equivalence. The result is a theorem about the actual flattened signal, not an unconnected pairwise calculation.
\excerpt{FiniteTransforms}{flat_lifting}
The companion direction, \code{lift_unlift_flat}, establishes reconstruction of arbitrary transformed coordinates as well. The finite Haar statements similarly include pair energy, finite energy, and a flat involution. These finite constructions are useful for matching array implementations, while the Hilbert files establish the infinite $\ell^2$ statements.

\subsection{A packet is an ordered composition}
At depth zero there is one empty word and its operator is the identity. A depth-$n+1$ word consists of a first branch and a remaining word. The first branch acts on the signal before the remaining packet.
\excerpt{WaveletPackets}{packet_recursion}
Here \code{Word n} is a recursively finite type. Each child map is a bounded continuous linear operator. The low and high norm bounds come from the two channel energy identity, so their adjoints are legitimate Hilbert adjoints.

\subsection{Induction preserves total energy}
The proof at depth $n+1$ first sums over the remaining word for each child. The induction hypothesis reduces that sum to the child's norm squared. Summing over the two children returns the parent's norm squared.
\excerpt{WaveletPackets}{packet_energy}
For depth two, this is the four-packet identity
\[
\|W_{00}x\|^2+\|W_{01}x\|^2+\|W_{10}x\|^2+\|W_{11}x\|^2=\|x\|^2.
\]
The theorem quantifies over every natural depth. Each individual depth has a finite packet family. It is not a theorem about convergence of an infinite depth packet expansion.

\subsection{Energy gives an operator reconstruction identity}
For finitely many bounded maps $W_i$ on the complex Hilbert space, an energy identity valid for every $x$ implies
\[
\sum_i W_i^*W_i=I.
\]
The formal proof uses inner product extensionality, the adjoint relation, and complex polarization through the library theorem it invokes.
\excerpt{WaveletPackets}{adjoint_identity}
This is a genuine adjoint identity. Writing a reconstruction map with a similar-looking formula would not suffice without showing that it is the adjoint of the analysis map.

Applying the resulting operator equality to a signal gives the source level synthesis theorem.
\excerpt{WaveletPackets}{packet_reconstruction}
The proof does not need every $W_i$ to be invertible. It needs the collective energy identity. The family can divide information among its channels while the sum of adjoint synthesis terms recovers the original signal exactly.

\Needspace{10\baselineskip}\section{Known edges and measurable perturbations}
\subsection{Exact retention separates the error algebraically}
Write $f=s+j$, where $s$ is the background and $j$ is a known jump packet. If the reconstruction is $Qs+j$, the retained jump cancels:
\[
f-(Qs+j)=s-Qs.
\]
\excerpt{KnownEdge}{known_edge}
The cancellation theorem needs only an additive commutative group and an arbitrary function $Q$. It does not require $Q$ to be linear. This is sufficient for cancellation because $Q$ is applied to $s$ before the same $j$ is restored. A separate claim that commutes $Q$ with addition, or estimates an operator commutator, would need the appropriate linearity hypotheses.

Taking norms yields exact equality of the full reconstruction error with the background error. A finite sum of known packets cancels in the same way. None of these statements estimates an unknown edge or proves that a data driven algorithm finds the correct mask.

\subsection{Estimated packets have two errors}
For $j=a\mathbf1_S$ and $\widehat j=b\mathbf1_T$,
\[
j-\widehat j=(a-b)\mathbf1_S+b(\mathbf1_S-\mathbf1_T).
\]
The first term changes the amplitude on the original mask. The second changes the mask while holding the estimated amplitude fixed. The symmetric difference $S\triangle T$ contains exactly the points where the two masks disagree.

For measurable sets and $1\le p<\infty$, the resulting bound is
\[
\|a\mathbf1_S-b\mathbf1_T\|_p
\le |a-b|\mu(S)^{1/p}+|b|\mu(S\triangle T)^{1/p}.
\]
The formal theorem lives first in the extended nonnegative reals, so infinite measures are handled without silently converting infinity to an ordinary real number.
\excerpt{MaskPerturbation}{mask_bound}
The notation \code{ENNReal} permits $\infty$, and \code{eLpNorm} is an extended seminorm. The hypotheses \code{hp : 1 ≤ p} and \code{hpTop : p ≠ ∞} govern the triangle inequality and finite-exponent formula. A finite real valued corollary requires the displayed terms to be finite, for example through finite mask and symmetric-difference measures.

\subsection{A concrete perturbation budget}
\begin{microexample}[title={An amplitude and mask budget}]
Suppose $p=2$, $\mu(S)=1$, $\mu(S\triangle T)=0.01$, $a=2$, and $b=1.9$. The bound is
\[
0.1\sqrt1+1.9\sqrt{0.01}=0.29.
\]
If the background approximation error is at most $0.04$, the full error is at most $0.33$ by the additional triangle inequality in \code{full_signal_perturbation}. The numbers illustrate a deterministic budget conditional on the amplitude, mask, and background estimates. A statistical procedure achieving those estimates requires additional work.

\end{microexample}

\subsection{Pointwise and essential supremum}
The source first proves a bound at every point, then transfers it to the essential supremum seminorm.
\excerpt{MaskPerturbation}{sup_bound}
The displayed bound uses whether $S\triangle T$ is empty as an actual set. An essential-supremum argument may ignore a null set, but a pointwise supremum cannot. This particular theorem remains a valid, sometimes coarse, bound when the sets differ only on a null set. Its statement should not be silently replaced by an assertion that pointwise and almost-everywhere equality coincide.

\subsection{The open obligations remain identifiable}
The stable edge-packet recovery conjecture asks for specified signal classes, contrast and separation assumptions, and a noise model under which a multiscale estimator controls both mask error and amplitude error at rates separated from smooth approximation. The formal perturbation bound explains how such estimates would enter reconstruction. It does not construct that estimator or prove those rates.

The edge-packet approximation-class conjecture asks for natural function spaces whose norms distinguish smooth approximation, boundary geometry, packet amplitude, and packet identification. The current declarations do not define and close that proposed theory. Infinite-tree consistency, minimax claims, and optimal Smith sampling also remain outside the twenty-result formal scope.

These boundaries matter when reading Part II. Its broader deterministic and geometric theorems can develop the mathematics without turning every contextual statement into a Lean-certified target. The final traditional record and the separate Part II crosswalk identify that distinction explicitly.

\Needspace{10\baselineskip}\section{Complete prunings, global selection, and stability}
\subsection{A finite tree carries positive finite branching}
The tree type permits a leaf or a branch with $n+1$ children. Thus every branch has at least one child, and unary branches are included.
\excerpt{FiniteAdaptive}{tree}
The tree itself is a finite inductive object. A cut is another typed inductive object: stop at the current root, or choose one cut in each child. The type prevents selecting an arbitrary collection that fails to cover a terminal path.
\excerpt{FiniteAdaptive}{cut}
A singleton leaf has only its stop cut. At a unary branch, stopping at the root and descending to its child are distinct cuts even if node labels happen to coincide. Formal nodes carry their positions in the tree, so equal labels do not collapse different nodes.

\subsection{Conservation is an explicit hypothesis}
For an additive quantity $p$ at each node, conservation says that the parent's value is the sum of the child-root values and that this continues recursively.
\excerpt{FiniteAdaptive}{conserves}
By induction on the cut, its sum equals the root value. The quantity may be a scalar, a vector, or an operator. The argument uses addition and the conservation identity. It does not derive conservation from an arbitrary tree topology.

For projector valued packets, one supplies both conservation and a root equal to the identity operator. Every complete pruning then reconstructs the signal. This makes selection compatible with exact synthesis: changing the cut changes the representation while preserving the total operator.

\subsection{From recursive cuts to covering antichains}
The manuscript describes a complete pruning as a node set that meets each root to terminal path exactly once. The source builds a positional node type, terminal paths, and a membership predicate for the nodes selected by a recursive cut. It proves the equivalence with the set based covering-antichain description.
\excerpt{CompletePrunings}{covering}
The file also proves that every complete set comes from a cut and that the cut is determined by its selected nodes. Together these results give \code{pruningEquiv}. This is why later theorems apply to every complete covering antichain, not just a preselected list of recursive examples.

\subsection{Literal sums over selected nodes}
\code{PruningSums} connects the recursive cut sum with the actual finite sum over the selected node set. Its operator theorem provides both the operator equality and the applied-signal equality.
\excerpt{PruningSums}{synthesis}
The hypotheses \code{h : conserves p t} and \code{hr : p (label t) = ...id...} are inputs. There is no assertion that an arbitrary family of packet operators indexed by a tree automatically satisfies them.

\subsection{Existence of a minimum allows nonadditive objectives}
The manuscript's general criterion has the form
\[
J(S)=R(S)+\eta L(S)+\tau K(S),
\]
where $R$, $L$, and $K$ can depend on the complete pruning as a whole. Finiteness alone ensures a minimum is attained.
\excerpt{PruningSums}{objective}
No positivity or additivity of these functions is needed for the finite attainment theorem. Such assumptions can be useful in applications, but the logical reason for attainment here is a nonempty finite search space.

The dynamic program below solves an additive criterion $\sum_{v\in S}c(v)$. The general nonadditive existence result does not imply that the same recurrence efficiently minimizes an arbitrary coupled objective. A user of the recurrence must verify the required additive decomposition.

\subsection{Strict gaps require competitors}
If $S_*$ is a strict minimizer and there is another pruning, the minimum positive objective difference is attained by a real competitor.
\excerpt{PruningSums}{gap}
The explicit \code{ho} hypothesis excludes the empty competitor family. For a singleton pruning space, the selected pruning is unchanged for every objective, so a separate singleton theorem is the natural statement. A tie between distinct prunings gives zero separation and does not satisfy a positive strict-gap hypothesis.

\subsection{The bottom up evaluator}
For an additive node cost, the optimal value obeys
\[
V(v)=\min\left(c(v),\sum_{u\text{ child of }v}V(u)\right).
\]
The selector stops at the node when the two alternatives tie. The formal evaluator returns the cut, its cost, and three operation counters together.
\excerpt{PruningAlgorithm}{solve}
The child results are introduced once as a family, and their stored costs are summed. The correctness theorem identifies both outputs with the specification: cost equals \code{value}, and cut equals \code{best}.
\excerpt{PruningAlgorithm}{solve_correct}
The proof uses tree induction. In a branch, the induction hypotheses identify all cached child values with their specification values. The same comparison then controls both the evaluator and the specification, including equality.

\subsection{An explicit small tree}
Consider a root $r$ with two leaf children $a$ and $b$. Its complete prunings are $\{r\}$ and $\{a,b\}$. Set $c(r)=7$, $c(a)=2$, and $c(b)=3$. The split costs $5$ and is selected. The gap to the other pruning is $2$.

Changing $c(r)$ to $5$ creates a tie. The selector now stops at the root by convention. This deterministic choice is defined, but it is not protected by a positive separation margin. An arbitrarily small perturbation can change which pruning has smaller cost. Tie handling and strict stability address different questions.

If $p(r)=p(a)+p(b)=I$, both prunings synthesize exactly. At a tie, two representations can exchange preference while both preserve the full signal through the same synthesis identity.

\subsection{What the linear-work count charges}
The source proves one visit per node, one comparison per internal node, and one addition per edge. The last identity is expressed as additions plus one equals the node count. Combining them yields
\excerpt{PruningAlgorithm}{linear_work}
Node costs are supplied inputs. Their estimation from data, the cost of exact real arithmetic, output serialization, and the runtime behavior of a particular compiler are not charged by this theorem. The claimed bound is precise for its stated unit operation model. It is not a wall clock performance certificate or a cost bound for a nonadditive objective.

\subsection{The exact worst pruning error envelope}
Suppose local perturbations obey $|d(v)-c(v)|\le e(v)$. The maximum possible sum of the error weights along a complete pruning has the recurrence
\[
E(v)=\max\left(e(v),\sum_{u\text{ child of }v}E(u)\right).
\]
\excerpt{FiniteAdaptive}{envelope}
The source proves both that every cut weight is at most this envelope and that some cut attains it. It is therefore the exact maximum of the additive weights over complete prunings, rather than a loose sum over all nodes.

For the two-leaf example, if $e(r)=0.3$ and $e(a)=e(b)=0.2$, then $E(r)=\max(0.3,0.4)=0.4$. Both candidate objective values are perturbed by at most $0.4$. The original gap is $2$, and $2E(r)=0.8<2$, so the selected split is protected.

\subsection{Why the margin contains a factor of two}
When comparing a competitor with the minimizer, the competitor can move down by the full error budget while the minimizer moves up by the same budget. Their separation can shrink by twice the budget. The antichain theorem preserves the selected pruning under the strict inequality.
\excerpt{PruningSums}{stability}
Equality at the boundary is insufficient in general. With an original gap $\gamma=2$ and budget $E=1$, the winner can move up by one and the competitor down by one, producing a tie. A deterministic tie rule can then change the selected object. The theorem's strict inequality is part of its content.

\subsection{Probability is transferred from a joint event}
The probabilistic statement assumes one event on which all required local inequalities hold simultaneously. Deterministic stability then makes that event a subset of the event that the selector is unchanged.
\excerpt{PruningAlgorithm}{probability}
Separate marginal statements for each node do not automatically give the same probability for their intersection. One needs a supplied joint event or an additional argument, such as a correctly budgeted union bound. No independence assumption or tail concentration theorem is inferred by the source.

The theorem is written for a measure with an explicit lower bound \code{1-δ ≤ μ good}. Lean's measure notation also evaluates arbitrary sets through the associated outer measure, so the set-inclusion proof does not itself establish event measurability. For ordinary success probabilities, use a probability measure, measurable random node costs, and a measurable joint good event. The selector on a finite tree is then determined by finitely many measurable comparisons, making its success event measurable. Finite tree structure alone cannot make arbitrary random cost functions measurable. The final traditional edition makes these probability hypotheses explicit. In Part II, consistency statements based only on probabilities tending to one are interpreted as convergence in probability. Almost-sure conclusions need additional hypotheses and an argument such as summable failure probabilities.

\Needspace{10\baselineskip}\section{Formal cyclic products and projective Smith actions}
\subsection{A formal residue projector filters coefficients}
For a formal power series $f=\sum_{n\ge0}f_nX^n$, the $j$th cyclic projector retains the degrees congruent to $j$ modulo $N$.
\excerpt{CyclicSeries}{cyclic_projector}
The ambient coefficient ring is a commutative semiring. There is no evaluation point and no convergence assumption in this definition. Multiplication is the formal Cauchy product, whose individual coefficients are finite sums.

The projector laws express idempotence, disjointness of different residue classes, and resolution of the identity for $N>0$. A pure packet in residue $i$ multiplied by one in residue $j$ lies in residue $i+j$ modulo $N$.

\subsection{The complete cyclic product law}
Decompose both factors into packets, distribute multiplication, and retain only pairs whose residue labels add to the requested output label.
\excerpt{CyclicSeries}{cyclic_product}
For $N=2$, write $f=f_0+f_1$ and $g=g_0+g_1$. Then
\[
(fg)_0=f_0g_0+f_1g_1,\qquad (fg)_1=f_0g_1+f_1g_0.
\]
The subscripts label degree parity. They are not independent scalar coefficients. This formula explains nonlinear coupling between packet labels without leaving the formal-series setting.

\subsection{Powers include the empty product}
The power theorem sums over all $q$-tuples of residue labels, retaining those whose sum has the requested residue.
\excerpt{CyclicSeries}{cyclic_power}
At $q=0$, the type \code{Fin 0 → Fin N} has one function, the empty assignment. Its sum of labels is zero and its product is one. Thus the formula returns the constant series one for the zero residue and zero otherwise. This boundary case is part of the statement, not an unstated convention added afterward.

The foundational coefficient-projector proofs identify their adaptation from Paper IV in source comments and in the cumulative census. Reusing them in a new application does not make them independent additional discoveries. Analytic root-of-unity formulas require their own convergence domain when formal series are evaluated.

\subsection{Use homogeneous coordinates at poles}
For a nonzero complex parameter $a$, the normalized homogeneous Smith action is
\[
A_a\binom uv=\frac12
\begin{pmatrix}1+a&1-a\\1-a&1+a\end{pmatrix}\binom uv.
\]
\excerpt{SmithAction}{smith_action}
The source proves $A_aA_b=A_{ab}$, $A_1=I$, and preservation of nonzero pairs when $a\ne0$. Multiplying both input coordinates by the same nonzero scalar multiplies both output coordinates by that scalar.

A projective point is a nonzero pair modulo this common scaling. The source constructs the equivalence relation explicitly.
\excerpt{SmithCyclic}{smith_quotient}
The quotient is the full projective line in this coordinate model. It includes $[0:1]$, representing zero impedance, and $[1:0]$, representing infinity. The action respects scaling, so it descends to the quotient.

\subsection{Cyclic composition}
The parameter $\omega_N=\exp(2\pi i/N)$ is nonzero and satisfies $\omega_N^N=1$ for $N>0$. The resulting projective actions compose by addition of exponents modulo $N$.
\excerpt{SmithCyclic}{smith_composition}
This source target establishes the modulo-$N$ composition law. It should not be described as a separate formal proof that the action has exact order $N$ unless the required primitivity and nontriviality converse are established as additional formal declarations. The broader traditional treatment contains further action analysis; its prose scope is separately recorded.

\subsection{Cayley coordinates explain the action}
In the pair coordinates $(u-v,u+v)$, the action multiplies the first coordinate by $a$ and leaves the second fixed.
\excerpt{SmithCyclic}{smith_cayley}
On a finite chart where the denominators are nonzero, this corresponds to rotation or multiplication of the reflection coefficient. The rational impedance formula is
\[
z\longmapsto\frac{(1+a)z+(1-a)}{(1-a)z+(1+a)}.
\]
A vanishing denominator means the image is the projective point at infinity. Lean's field division is totalized, so a bare identity of division expressions does not by itself assert that a displayed chart exists. The homogeneous action carries the geometry through such poles, while chart composition theorems carry explicit denominator conditions.

\subsection{The half-turn swaps zero and infinity}
For $N=2$, $\omega_2=-1$ and the homogeneous action swaps the two coordinates.
\excerpt{SmithCyclic}{smith_swap}
On the finite nonzero chart this gives $z\mapsto z^{-1}$. At $z=0$ it maps $[0:1]$ to $[1:0]$. At infinity it does the reverse. The reciprocal formula is therefore a chart description of a globally defined projective transformation, not a rule that discards poles.

\subsection{Dyadic and odd masks factor by CRT}
If $N=2^aM$ with $M$ odd, the two factors are coprime. An integer satisfies a residue condition modulo their product precisely when it satisfies both factor conditions. The source proves the corresponding equality of integer-residue masks.
\excerpt{CRT}{crt}
This is an exact mask factorization. It does not prove that a particular nonuniform sample allocation is optimal, nor does it provide a minimax sampling theorem. It supplies the arithmetic structure on which more specialized sampling constructions can be built.

\Needspace{10\baselineskip}\section{Trust boundaries and reproduction}
\label{sec:reproduce}
\subsection{What the inherited audit actually checks}
The preserved compiler log records 225 explicit theorem targets and 827 namespace declarations. The audit driver first requires each target to be present as a theorem. It then traverses the \code{KomorebiPaperV} namespace, rejects unsafe declarations, rejects local axiom and opaque declarations, and collects transitive axiom dependencies.

Only \code{propext}, \code{Classical.choice}, and \code{Quot.sound} are permitted. The audit checks the dependencies of the compiled declarations. It does not prove that English prose is a correct paraphrase of those declarations, that a Python implementation matches a Lean definition, or that an application has supplied a theorem's hypotheses.

The current guide build reconciles each target name with its immutable source location and inherited axiom-output record. It verifies pinned source hashes and performs a lexical scan for proof placeholders and prohibited source constructs. These checks protect provenance and locate obvious source-policy violations. They do not replace elaboration, type checking, or a new kernel traversal.

\subsection{The two dates of evidence}
The release's raw kernel verification timestamp is 6 September 2026 at approximately 03:18 UTC. This monograph companion is dated 9 October 2026. The compiler transcript describes the preserved v1.1.0 sources and their pinned dependencies. Its date remains visible beside the current checks.

This edition freshly ran the twenty exact model tests, comprising 305,933 counted assertions, and all sixteen verifier control methods. It ran the 33 current companion pytest tests and the 27 overlapping historical tests with no skips. These are finite tests of implementations and models, with separate contracts for their tested scope.

All 52 stored floating numerical metrics and eighteen discrete result leaves were reproduced exactly in the available environment. Eighteen companion PDF/PNG figure pairs were regenerated. The broader Annals companion separately passed 11,052 mixed exact and numerical checks across thirteen groups. Six bounded computational manifests were validated against their inputs. The symbolic visual suite passed 181 checks and rejected seven false controls. The extension examples separately passed 38 exact checks and rejected four false controls.

These counts refer to the recorded tests and their declared ranges. Overlapping suites are retained as provenance and are not combined into a theorem denominator. In particular, floating computations and bounded models do not extend the twenty-result formal contract or close either edge recovery conjecture.

No fresh Lean process ran because \code{lean}, \code{lake}, and \code{elan} were unavailable. Current static reconciliation checks the actual source bytes, 225 declaration locations, 225 audit directives, nine dependency revisions, and 225 saved axiom records. It does not run Lean's elaborator or kernel. The complete package's current computational review and per-run manifests record the fresh tests and their limits.

The recovered October 5 edition lacks seventeen verification artifacts, including its \code{verify_available.py} runner and some report and figure files. The surviving formal sources, September compiler logs, and saved axiom audit are preserved. Current scoped replays have their own manifests and adapters. The October 5 available-checks gate was not rerun or authenticated from a complete historical inventory. The inventory appears in \code{provenance/HISTORICAL_RECOVERY_GAPS.json}. Missing historical evidence is not recreated by a current passing test.

\subsection{Reproduce from a working copy}
Use a writable copy of the original release to run its verification entry points, because they write status and log files beside the sources. Keep the packaged original snapshot unchanged as the provenance reference. The following commands start in that copied release root, which contains \code{VERIFY_PAPER.py}, \code{formalization}, \code{companion}, and \code{verification}.

The retained commands below describe the September formal release entry points. They do not reproduce the missing October 5 available-checks runner. The original formal-release pipeline expects the pinned Lean toolchain, its pinned dependencies, the Python requirements in \code{companion/requirements.txt}, \code{pypdf}, pdfLaTeX, and Poppler tools. Once those prerequisites are available, obtain the mathlib cache in the Lean project:
\begin{Verbatim}[fontsize=\small,frame=leftline,breaklines=true]
cd formalization/lean
lake exe cache get
cd ../..
python3 VERIFY_PAPER.py
\end{Verbatim}
This full entry point checks the reviewed source hashes, calls the Lean verifier, runs the original companion tests, repeats the optimized static checks, replays numerical experiments, regenerates figure pairs, and rebuilds the selected 39-page source manuscript. It reproduces the original integrated formal-core release, not the later combined two-part traditional PDF.

For the formal source and exact-model subset, use:
\begin{Verbatim}[fontsize=\small,frame=leftline,breaklines=true]
python3 verification/verify.py
\end{Verbatim}
This command checks the actual toolchain version, builds the local Lean library, checks pinned dependency revisions and tracked-file cleanliness, and runs \code{KernelAudit.lean} with warnings treated as errors. If Lake is unavailable, the verifier reports \code{BLOCKED_NO_LAKE} and exits with code two. It does not turn passing Python tests into a successful Lean result.

For a deliberately static and exact-model run with no Lean claim:
\begin{Verbatim}[fontsize=\small,frame=leftline,breaklines=true]
python3 -O verification/verify.py --static-only
\end{Verbatim}
The \code{STATIC_AND_EXACT_ONLY} status and \code{lean_kernel_checked=false} are expected for that route. Python optimization is an intentional verifier control; the program's critical requirements are explicit checks rather than removable Python \code{assert} statements.

\subsection{Manual inspection of the formal build}
To inspect the principal operations separately, work inside \code{formalization/lean}:
\begin{Verbatim}[fontsize=\small,frame=leftline,breaklines=true]
lake env lean --version
lake build KomorebiPaperV
lake env lean -DwarningAsError=true KernelAudit.lean
\end{Verbatim}
A successful build without the audit is incomplete evidence for the release's declared axiom policy. A successful audit of changed sources is a new result that should be recorded with new source hashes. Do not reuse an old transcript as though it described a modified source tree.

The historical audit's final success line is:
\begin{Verbatim}[fontsize=\small,frame=leftline,breaklines=true]
PAPER_V_CANDIDATE_TARGET_AUDIT_PASS: 225 theorem targets; 827 namespace declarations
\end{Verbatim}
The whole transcript, compiler version, source pins, and exit status matter. Searching for this line in a file is provenance reconciliation, not a substitute for executing the audit.

\subsection{The historical runtime adapter}
The original environment used a narrow executable self-location adapter. Its ordinary process-specific executable path returned a permission error, while \code{/proc/self/exe} was accessible. The source adapter redirected only the current process's own executable-link query. The bundled documentation records that the distributed compiler and kernel were unchanged.

The retained positive control accepted a valid arithmetic theorem and the negative control rejected an invalid proof with exit code one. These are useful runtime controls, but the full source build and namespace audit remain the decisive historical evidence. The adapter is environment-specific and unnecessary on a normal installation. Toolchain binaries and compiled dependency caches are not included in this package.

\subsection{Rebuild this guide}
The guide itself needs XeLaTeX and the TeX packages named in \code{preamble.tex} and \code{komorebi-monograph.sty}. Its P052, Nimbus Sans, DejaVu Mono, and Latin Modern fonts must be installed externally. The builder resolves them with \code{kpsewhich}, verifies the staged bytes, and records their hashes. It never installs dependencies or bundles font binaries. From its editable directory:
\begin{Verbatim}[fontsize=\small,frame=leftline,breaklines=true]
python3 build.py
\end{Verbatim}
This copies the complete authored guide tree to an isolated directory and compiles the supplied TeX and generated catalog files in four passes. It executes no Lean code and does not regenerate mathematics. The separate \code{build_guide_data.py} script can regenerate the catalogs when supplied the immutable release root and the traditional paper's final crosswalk. It verifies hashes and rejects mismatched contracts before writing the guide's data products.

The raw excerpt files preserve source text. Their rendering files escape TeX-special characters and route rare mathematical glyphs to a suitable font. The PDF is a reading surface. The full source files retain the actual namespace and typeclass context needed for Lean compilation.

\subsection{A theorem-use checklist}
Before using a result, identify its ambient type and explicit hypotheses, locate the exact declaration, and check how the application supplies each hypothesis. For an FIR bank, confirm finite Laurent support and delay units. For a transform equality, supply weighted summability. For pruning synthesis, supply conservation. For selector stability, establish the strict gap and joint local-error bound. For a Smith chart, check denominator nonvanishing or use the projective action.

This is the recurring pattern of the formal development: construct the object in the correct space, prove the bridge to the intended model, then use the identity whose hypotheses have actually been supplied.

\clearpage
\section{The complete twenty-result correspondence}
\label{sec:ledger}
Each entry gives the scoped formal statement, its principal declarations, and the traditional manuscript locator. The 225-target catalog includes the supporting theorems as well. Principal targets may recur because one synthesis fact supports more than one manuscript statement. These are navigation relationships, not additive theorem counts. Traditional locators use printed Arabic pages. The unnumbered cover and Roman opening pages have separate physical PDF indices, preserved in the companion crosswalk.
\input{generated/result_ledger}
\clearpage
\section{The broader traditional developments}
\label{sec:extension}
The final traditional edition preserves two substantial source strands. Part I contains the twenty-result formal core. Part II contains sixty-three proof records from the broader Annals development. Together they produce 83 proof records, with overlap in mathematical content. The number 83 counts presented proofs. It does not replace the twenty-result Lean denominator or establish 83 independent new results.

The table below connects groups of Part II developments to nearby formal-core concepts. ``Nearby'' is a reading aid. It does not assert a formal reduction or that the stated Part II result has been compiled in Lean.
\begin{longtable}{>{\raggedright\arraybackslash}p{.27\textwidth}>{\raggedright\arraybackslash}p{.29\textwidth}>{\raggedright\arraybackslash}p{.33\textwidth}}
\toprule Part II topic & Nearest formal core & Additional boundary \\ \midrule\endhead
Operator partitions and refinement & Results 1.1--2.3, \code{HilbertPackets} & Broader operator settings or frames require their own assumptions and proofs. \\ \addlinespace
Polyphase symbols and banks & Results 3.1--4.3, \code{FilterBank}, \code{ZTransform} & Unit-circle and $L^2$ boundary interpretations must be distinguished from off-circle weighted summability. \\ \addlinespace
Frame trees and adaptive objectives & Results 11.2--11.5, \code{PruningSums} & Local frame bounds are nonnegative. The dynamic-programming result applies to additive costs. \\ \addlinespace
Known edges, masks, and commutators & Results 6.1--7.1, \code{KnownEdge}, \code{MaskPerturbation} & A commutator using $Q$ as an operator requires bounded linearity. Cancellation alone does not. \\ \addlinespace
Boundary tubes and edge geometry & Mask perturbation theorem 7.1 & Tube estimates use their local range $\delta<\delta_0$. Geometric constants and regularity are additional data. \\ \addlinespace
Continuous Haar edge models & Discrete Haar results 4.2 and 5.1 & Half-line step functions need care: their wavelet coefficients can be used even when the steps themselves are not in $L^2(\mathbb R)$. \\ \addlinespace
Random or asymptotic selection & Finite joint-event transfer 11.5 & Probability bounds tending to one establish convergence in probability. Almost sure claims need stronger evidence. \\ \addlinespace
Nonlinear cyclic products & Result 12.1, \code{CyclicSeries} & Analytic evaluation and stronger functional-analytic claims are additional to formal coefficient identities. \\ \addlinespace
Smith geometry and sampling & Results 13.2--14.1, \code{SmithCyclic}, \code{CRT} & The formal composition law alone is not a target proving exact-order converse or optimal sampling. \\ \bottomrule
\end{longtable}
The final traditional audit records the precise scope repairs applied to Part II. This guide preserves the resulting hypotheses in its discussion. It does not incorporate altered claims into unchanged Lean targets. The following inventory provides a locator for every Part II proof record and explicitly keeps its formal status separate.

\input{generated/extension_crosswalk}
\clearpage
\section{All 225 theorem targets}
\label{sec:targets}
This catalog reproduces every qualified target from the frozen \code{candidate_contract.json}, grouped by its actual source module. Line numbers refer to the immutable bundled Lean sources. Every target was reconciled with its declaration line, audit directive, and inherited axiom-output record. This reconciliation does not perform Lean elaboration.
\input{generated/target_catalog}
\clearpage
\section{Source architecture and exact dependency pins}
\label{sec:imports}
The import list below is extracted from the 25 Lean files: the 23 proof modules, the aggregate import file, and the audit driver. Imports beginning with \code{Mathlib} are external mathematical dependencies. Imports beginning with \code{KomorebiPaperV} are local. The import graph describes compilation dependencies, which can be broader than the logical dependency of an individual theorem.
\input{generated/imports}
\subsection{Pinned external revisions}
The package uses Lean \code{leanprover/lean4:v4.34.0-rc2}. The following revisions are copied from \code{lake-manifest.json}. They specify the historical dependency state. This guide preserves those revisions.
\input{generated/pins}
\end{document}
